\section*{Capítulo \ref{ch:b3:topology} --- Topologia geral}

\begin{solution}{exo:b3:topology:1}
(a) Quatro \omterm{def:b3:topology:topology}{topologias} sobre $\{a,b\}$: a indiscreta
$\{\varnothing, X\}$; a discreta; as duas \omterm{def:b3:topology:topology}{topologias} de \emph{Sierpiński}
$\{\varnothing, \{a\}, X\}$ e $\{\varnothing,
\{b\}, X\}$. As duas últimas são \omterm{def:b3:topology:continuity}{homeomorfas} (troque $a, b$): três
classes. \omterm{def:b3:topology:connected}{Conexas}: todas, exceto a discreta (só a
\omterm{def:b3:topology:topology}{topologia} discreta contém um aberto-fechado próprio não vazio).
Hausdorff: só a discreta (nas outras, a única \omterm{def:b3:topology:topology}{vizinhança} de
$b$ é $X$, ou a de $a$).

(b) $\Q$: \omterm{def:b3:topology:interior}{interior} $\varnothing$ (todo intervalo contém
irracionais), \omterm{def:b3:topology:interior}{fecho} $\R$ (densidade), fronteira $\R$. $[0,1)
\cup \{2\}$: \omterm{def:b3:topology:interior}{interior} $(0,1)$, \omterm{def:b3:topology:interior}{fecho} $[0,1]\cup\{2\}$,
fronteira $\{0, 1, 2\}$. $\{1/n\}$: \omterm{def:b3:topology:interior}{interior} $\varnothing$,
\omterm{def:b3:topology:interior}{fecho} $\{1/n\}\cup\{0\}$, fronteira o próprio \omterm{def:b3:topology:interior}{fecho}.
\end{solution}

\begin{solution}{exo:b3:topology:2}
(a) Todo aberto $V \subseteq Y$ é uma união $\bigcup B_j$ de conjuntos
da \omterm{def:b3:topology:basis}{base}, e $f^{-1}(V) = \bigcup f^{-1}(B_j)$: se estes últimos são
abertos, o primeiro também é; a recíproca é trivial.

(b) $f^{-1}\bigl(\intoo{\frac12}{\frac32}\bigr) = [0, \infty)$
não é aberto: $f$ não é \omterm{def:b3:topology:continuity}{contínua}; a \omterm{def:b3:topology:continuity}{continuidade} à direita em cada
ponto é clara ($f$ é localmente constante à direita de cada
ponto). Os conjuntos $[a, b)$ satisfazem o critério de \omterm{def:b3:topology:basis}{base}
(\cref{def:b3:topology:basis}): $[a,b)\cap[c,d) = [\max(a,c),
\min(b,d))$. A \omterm{def:b3:topology:topology}{topologia} de Sorgenfrey é mais fina que a
usual, pois $(a, b) = \bigcup_{a < c < b}[c, b)$; estritamente:
$[0,1)$ é aberto de Sorgenfrey, mas não é aberto usual. \omterm{def:b3:topology:continuity}{Continuidade} a partir da
reta de Sorgenfrey em $x$: uma \omterm{def:b3:topology:topology}{vizinhança} de Sorgenfrey de $x$
contém um conjunto da \omterm{def:b3:topology:basis}{base} $[a, b) \ni x$, logo contém $[x, b)$ ---
de modo que a condição de \omterm{def:b3:topology:continuity}{continuidade} se lê: para todo $\varepsilon > 0$
existe $\delta > 0$ com $\abs{f(t) - f(x)} < \varepsilon$
sempre que $x \leq t < x + \delta$. Isso é exatamente a
\omterm{def:b3:topology:continuity}{continuidade} à direita em $x$.
\end{solution}

\begin{solution}{exo:b3:topology:3}
Dois abertos não vazios têm complementares finitos, de modo que sua
interseção tem complementar finito: não vazia (o conjunto é
infinito) --- dois pontos nunca têm \omterm{def:b3:topology:topology}{vizinhanças} disjuntas: não é
Hausdorff, logo não é metrizável
(\cref{def:b3:topology:hausdorff}). Seja $(x_n)$ injetiva e
$x$ arbitrário: uma \omterm{def:b3:topology:topology}{vizinhança} de $x$ contém um aberto cofinito
$U$; os finitos pontos de $X\setminus U$ são atingidos por, no
máximo, finitos índices (injetividade), de modo que uma cauda da
sequência está em $U$: $x_n \to x$, para \emph{todo} $x$. A
demonstração da unicidade precisa de duas \omterm{def:b3:topology:topology}{vizinhanças} disjuntas para separar
dois supostos limites --- precisamente o que falha aqui.
\end{solution}

\begin{solution}{exo:b3:topology:4}
(a) A \omterm{def:b3:topology:continuity}{continuidade} é por construção
(\cref{def:b3:topology:constructions}). Abertura: um aberto $W
\subseteq X \times Y$ é uma união de blocos $U_i\times V_i$, e
$\pi_X(W) = \bigcup U_i$ (blocos não vazios se projetam sobre seus
fatores), aberto. Não fechada: $H = \{(x,y) : xy = 1\}$ é fechado
em $\R^2$ (pré-imagem de $\{1\}$ pelo produto, \omterm{def:b3:topology:continuity}{contínuo}),
mas $\pi_X(H) = \R\setminus\{0\}$ não é fechado.

(b) ($\Rightarrow$) As projeções são \omterm{def:b3:topology:continuity}{contínuas}. ($\Leftarrow$)
Seja $x^{(k)} \to x$ componente a componente e $W = \prod U_n$ uma \omterm{def:b3:topology:topology}{vizinhança}
básica de $x$, com $U_n = X_n$ salvo para $n \in F$
finito. Para cada $n \in F$, $x_n^{(k)} \in U_n$ para $k \geq
K_n$; para $k \geq \max_F K_n$, $x^{(k)} \in W$. A fórmula $d$
define uma métrica (cada parcela é uma, a menos da verificação
usual de que $\min(1, d_n)$ o é); suas bolas: $B_d(x,
2^{-N})$ contém o bloco básico $\{y : d_n(x_n,y_n) < 2^{-N},\
n \leq N'\}$ para $N'$ adequado (a cauda $\sum_{n > N'} 2^{-n}$
é pequena) e, reciprocamente, todo bloco básico contém uma $d$-bola:
as duas \omterm{def:b3:topology:topology}{topologias} têm as mesmas \omterm{def:b3:topology:topology}{vizinhanças} de cada ponto.

(c) Na \omterm{def:b3:topology:constructions}{topologia produto}, $x^{(k)} \to 0$ por (b). O conjunto
$U = \prod_n \intoo{-\tfrac1n}{\tfrac1n}$, aberto na \omterm{def:b3:topology:topology}{topologia} das caixas, contém $0$, mas
$x^{(k)} \notin U$ para todo $k$ (a $n$-ésima coordenada $1/k
\geq 1/n$ quando $n \geq k$): nenhuma cauda entra em $U$. A \omterm{def:b3:topology:topology}{topologia} das caixas
é estritamente mais fina e não é, para fins de convergência, uma
\omterm{def:b3:topology:topology}{topologia} do tipo produto.
\end{solution}

\begin{solution}{exo:b3:topology:5}
$\varphi \colon \R \to S^1$, $x \mapsto \eu^{2\iu\pi x}$, é
\omterm{def:b3:topology:continuity}{contínua}, sobrejetora e constante nas classes módulo $\Z$: ela
induz uma bijeção \omterm{def:b3:topology:continuity}{contínua} $\bar\varphi \colon \R/\Z \to
S^1$ (\cref{prop:b3:topology:universal}(b)). A projeção
$\pi$ é \emph{aberta}: para $U \subseteq \R$ aberto,
$\pi^{-1}(\pi(U)) = \bigcup_{n}(U + n)$ é aberto, de modo que $\pi(U)$ é
aberto. Hausdorff: sejam $\bar x \neq \bar y$; escolha
representantes com $\delta = \min_{n}\abs{x - y - n} > 0$;
as imagens dos intervalos de raio $\delta/3$ em torno de $x$ e
$y$ são abertas (abertura de $\pi$), contêm $\bar x, \bar y$ e
são disjuntas (dois pontos das pré-imagens distariam $< 2\delta/3$
módulo $\Z$). \omterm{def:b3:topology:compact}{Compacto}: $\R/\Z = \pi([0,1])$, imagem \omterm{def:b3:topology:continuity}{contínua} de
um \omterm{def:b3:topology:compact}{espaço compacto} (\cref{thm:b3:topology:compactprops}(2)). Agora
$\bar\varphi$ é uma bijeção \omterm{def:b3:topology:continuity}{contínua} de um \omterm{def:b3:topology:compact}{espaço compacto} no
\omterm{def:b3:topology:hausdorff}{espaço de Hausdorff} $S^1$: um \omterm{def:b3:topology:continuity}{homeomorfismo}
(\cref{cor:b3:topology:compacthomeo}).

Para (iii): a composta $[0,1] \hookrightarrow \R
\xrightarrow{\pi} \R/\Z$ é \omterm{def:b3:topology:continuity}{contínua}, sobrejetora e
identifica exatamente $0 \sim 1$: ela induz uma bijeção \omterm{def:b3:topology:continuity}{contínua}
$[0,1]/(0\sim1) \to \R/\Z$ de um \omterm{def:b3:topology:compact}{espaço compacto} (imagem
\omterm{def:b3:topology:continuity}{contínua} de $[0,1]$ pela projeção quociente) em um \omterm{def:b3:topology:hausdorff}{espaço de Hausdorff}:
um \omterm{def:b3:topology:continuity}{homeomorfismo}. Compondo: os três espaços são
\omterm{def:b3:topology:continuity}{homeomorfos}.
\end{solution}

\begin{solution}{exo:b3:topology:6}
(a) Uma cobertura de $K_1 \cup \dots \cup K_m$ se restringe a uma cobertura de
cada $K_j$: finitos abertos por peça bastam. Uma
interseção $\bigcap K_i$ é fechada no \omterm{def:b3:topology:compact}{compacto} $K_{i_0}$
(\omterm{def:b3:topology:compact}{compactos} são fechados no espaço métrico ambiente), logo
compacta.

(b) $x \mapsto d(x, F)$ é \omterm{def:b3:topology:continuity}{contínua} ($1$-lipschitziana) e
estritamente positiva em $K$ ($d(x, F) = 0$ significa $x \in \bar F =
F$); no \omterm{def:b3:topology:compact}{compacto} $K$ ela atinge um mínimo $m > 0$: $d(K, F)
\geq m$. Contraexemplo sem \omterm{def:b3:topology:compact}{compacidade}: $F_1 = \{n : n \geq
2\}$ e $F_2 = \{n + \frac1n : n \geq 2\}$ são fechados,
disjuntos, à distância $\inf_n \frac1n = 0$.

(c) Se $X$ é \omterm{def:b3:topology:compact}{compacto}, toda $f$ \omterm{def:b3:topology:continuity}{contínua} é limitada
(\cref{thm:b3:topology:compactprops}(2)). Reciprocamente, se $X$ não é
\omterm{def:b3:topology:compact}{compacto}, tome $(x_n)$ sem subsequência convergente
(\cref{thm:b3:topology:metriccompact}); passando a uma
subsequência, podemos supor os $x_n$ dois a dois distintos, e nenhum
ponto de $X$ é valor de aderência da sequência, de modo que
\[
r_n = \min\Bigl(2^{-n},\ \tfrac13\,d\bigl(x_n, \{x_m : m \neq
n\}\bigr)\Bigr) > 0,
\]
e as bolas $B(x_n, r_n)$ são duas a duas disjuntas (um ponto comum
à $n$-ésima e à $m$-ésima daria $d(x_n, x_m) < r_n +
r_m \leq \frac23 d(x_n, x_m)$). Defina
\[
f(x) = \sum_{n} n\,\max\Bigl(0,\ 1 - \frac{d(x,
x_n)}{r_n}\Bigr):
\]
em cada $x$ há no máximo uma parcela não nula, e $f(x_n) = n$.
\omterm{def:b3:topology:continuity}{Continuidade} em $x$: seja $y_k \to x$; cada valor $f(y_k)$ é
$0$ ou um único valor de pico $\mathrm{bump}_{n_k}(y_k)$.
Se algum índice $n$ ocorre uma infinidade de vezes, ao longo dessa
subsequência $f(y_k) = \mathrm{bump}_n(y_k) \to
\mathrm{bump}_n(x)$, que é igual a $f(x)$ (se $x \in B(x_n,
r_n)$, toda a cauda está nessa bola aberta e nenhum outro
índice ocorre; se $x \notin B(x_n, r_n)$, então
$\mathrm{bump}_n(x) = 0 = f(x)$, pois $x$, aderente à
$n$-ésima bola, não está em nenhuma outra bola aberta). Se $n_k \to \infty$:
$d(y_k, x_{n_k}) < r_{n_k} \leq 2^{-n_k} \to 0$, de modo que $x_{n_k}
\to x$, tornando $x$ valor de aderência de $(x_n)$ --- excluído; assim
esse caso diz respeito apenas a finitos $k$. Em todos os casos
$f(y_k) \to f(x)$: $f$ é \omterm{def:b3:topology:continuity}{contínua} e ilimitada.
\end{solution}

\begin{solution}{exo:b3:topology:7}
Suponha que nenhum $\delta$ funcione: para cada $n$ existe $A_n$ de
diâmetro $< 1/n$ que não está contido em nenhum $U_i$; tome $x_n \in
A_n$. Por \omterm{def:b3:topology:compact}{compacidade} (\cref{thm:b3:topology:metriccompact}),
uma subsequência $x_{n_k} \to x$; escolha $i$ e $r > 0$ com $B(x,
r) \subseteq U_i$. Para $k$ grande: $d(x_{n_k}, x) < r/2$ e
$1/n_k < r/2$, de modo que $A_{n_k} \subseteq B(x_{n_k}, 1/n_k)
\subseteq B(x, r) \subseteq U_i$ --- contradição. Heine:
dado $\varepsilon$, cubra $Y$ por bolas $B(y, \varepsilon/2)$;
as pré-imagens formam uma cobertura aberta de $X$; seja $\delta$ um
número de Lebesgue: se $d(x, x') < \delta$, o par $\{x, x'\}$
tem diâmetro $< \delta$, está contido em uma pré-imagem, e $d(f(x),
f(x')) < \varepsilon$.
\end{solution}

\begin{solution}{exo:b3:topology:8}
(a) $GL_n(\R) = \det^{-1}(\R\setminus\{0\})$ é aberto ($\det$ é
polinomial, logo \omterm{def:b3:topology:continuity}{contínuo}). Desconexo: $\det$ o leva sobre
$\R\setminus\{0\}$, e um \omterm{def:b3:topology:connected}{espaço conexo} tem imagens
\omterm{def:b3:topology:continuity}{contínuas} \omterm{def:b3:topology:connected}{conexas} (\cref{thm:b3:topology:connectedbasics}(2));
$\R\setminus\{0\}$ não é um intervalo. $GL_n(\C)$: para $A, B$
inversíveis, $p(\lambda) = \det\bigl((1-\lambda)A + \lambda
B\bigr)$ é um polinômio em $\lambda$ com $p(0), p(1) \neq 0$,
logo não identicamente nulo: tem finitas raízes em $\C$.
O plano menos finitos pontos é \omterm{def:b3:topology:pathconnected}{conexo por caminhos} (contorne
os pontos), de modo que existe um \omterm{def:b3:topology:pathconnected}{caminho} $\lambda(t)$
de $0$ a $1$ com $p(\lambda(t)) \neq 0$: $t \mapsto
(1-\lambda(t))A + \lambda(t)B$ é um \omterm{def:b3:topology:pathconnected}{caminho} em $GL_n(\C)$.

(b) $\det(A + \varepsilon I) = \chi_A(-\varepsilon)\cdot(\pm1)$
se anula para no máximo $n$ valores de $\varepsilon$: matrizes
inversíveis $A + \varepsilon I \to A$ quando $\varepsilon \to 0$:
densidade.
\end{solution}

\begin{solution}{exo:b3:topology:9}
(a) Seja $A \subseteq \Q$ contendo $p < q$ e tome um irracional
$\alpha \in (p, q)$: $A = (A \cap (-\infty, \alpha)) \sqcup
(A\cap(\alpha, +\infty))$ reparte $A$ em duas peças não vazias,
relativamente abertas: $A$ é desconexo. As componentes são os pontos.
A \omterm{def:b3:topology:compact}{compacidade} local falha: uma \omterm{def:b3:topology:topology}{vizinhança} compacta $V$ de $0$ em
$\Q$ contém $[-r, r]\cap\Q$ para algum $r > 0$, que é fechado
em $V$, logo \omterm{def:b3:topology:compact}{compacto}; mas uma sequência de racionais em $[-r,r]$
que converge (em $\R$) a um irracional não tem subsequência
convergente \emph{em $\Q$}: contradição com o~\cref{thm:b3:topology:metriccompact}.

(b) Seja $\gamma = (\gamma_1, \gamma_2) \colon [0,1] \to S$ um
\omterm{def:b3:topology:pathconnected}{caminho} com $\gamma(0) = (0,0)$ e $\gamma(1) \in \Gamma$. O
conjunto $\{t : \gamma_1(t) = 0\}$ é fechado e não contém
$1$; seja $t^*$ seu supremo, de modo que $\gamma_1(t^*) = 0$ e
$\gamma_1 > 0$ em $(t^*, 1]$. Pela \omterm{def:b3:topology:continuity}{continuidade} de $\gamma_2$ em
$t^*$, escolha $\delta$ com $\abs{\gamma_2(t) -
\gamma_2(t^*)} < \tfrac12$ para $t \in [t^*, t^* + \delta]$. Fixe
$t_0 = t^* + \delta$ (podemos supor $\leq 1$): $\gamma_1(t_0) > 0$,
e $\gamma_1$ assume todos os valores de $(0, \gamma_1(t_0)]$ em
$[t^*, t_0]$ (teorema do valor intermediário,
\cref{thm:b3:topology:connectedbasics}). Escolha $k$ grande com
$a_k = \frac1{\pi/2 + 2k\pi} < \gamma_1(t_0)$ e $b_k =
\frac1{3\pi/2 + 2k\pi} < \gamma_1(t_0)$: existem $s, s' \in
(t^*, t_0]$ com $\gamma_1(s) = a_k$, $\gamma_1(s') = b_k$;
como os pontos estão no gráfico, $\gamma_2(s) = \sin(1/a_k) =
1$ e $\gamma_2(s') = -1$. Ambos não podem estar a menos de $\frac12$ de
$\gamma_2(t^*)$: contradição. $S$ é \omterm{def:b3:topology:connected}{conexo}, mas não é
\omterm{def:b3:topology:pathconnected}{conexo por caminhos}.
\end{solution}

\begin{solution}{exo:b3:topology:10}
(a) $C_n$ consiste exatamente nos $x$ que admitem uma expansão
ternária com dígitos em $\{0, 2\}$ até a ordem $n$ (indução:
retirar os terços médios elimina o primeiro dígito $1$, etc.; os extremos
têm duas expansões, uma delas evitando os $1$), de modo que $C = \bigcap C_n$
é o conjunto das somas $\sum a_n3^{-n}$, $a_n \in \{0,2\}$. \omterm{def:b3:topology:compact}{Compacto}:
cada $C_n$ é uma união finita de intervalos fechados; $C$ é fechado
em $[0,1]$. \omterm{def:b3:topology:interior}{Interior} vazio: $C \subseteq C_n$, união de
intervalos de comprimento $3^{-n}$; um intervalo \omterm{def:b3:topology:interior}{interior} de comprimento
$\ell > 0$ caberia em um deles para todo $n$. \omterm{prop:b3:galois:perfect}{Perfeito}: dados
$x \in C$ e $m$, troque o dígito $a_m$: o novo ponto está em
$C$, é distinto e dista menos de $2\cdot3^{-m}$ de $x$.

(b) Se $x \neq y \in C$, suas expansões diferem pela primeira vez em alguma
ordem $k$; entre eles está um terço médio retirado
(a lacuna de ordem $k$ que separa o dígito $0$ do dígito $2$),
fornecendo um ponto $c \notin C$, $x < c < y$ (digamos): $C =
(C \cap (-\infty, c)) \sqcup (C \cap (c, \infty))$ reparte qualquer
subconjunto que contenha ambos os pontos. As componentes são os pontos.

(c) A aplicação dos dígitos $\beta\colon x \mapsto (a_n(x))$ é uma
bijeção sobre $\{0,2\}^\N$ (existência e unicidade das
expansões em $\{0,2\}$: sequências de dígitos distintas dão pontos à
distância $\geq 3^{-k} > 0$ na ordem em que diferem pela primeira vez, como no~\cref{pb:b3:topology:1}, questão 1). Ela é \omterm{def:b3:topology:continuity}{contínua}:
$\abs{x - y} < 3^{-m}$ força a coincidência dos $m$ primeiros dígitos
(mesmo cálculo), de modo que $\beta$ leva bolas pequenas em blocos
da \omterm{def:b3:topology:basis}{base}. Uma bijeção \omterm{def:b3:topology:continuity}{contínua} do \omterm{def:b3:topology:compact}{compacto} $C$ no
produto \omterm{def:b3:topology:hausdorff}{de Hausdorff} é um \omterm{def:b3:topology:continuity}{homeomorfismo}
(\cref{cor:b3:topology:compacthomeo}; o produto é \omterm{def:b3:topology:hausdorff}{de Hausdorff}:
separe numa coordenada em que difiram). Não enumerabilidade: diagonal de
Cantor em $\{0,2\}^\N$. Por fim, $\{0,2\}^\N \times
\{0,2\}^\N \cong \{0,2\}^\N$ intercalando dígitos (um
\omterm{def:b3:topology:continuity}{homeomorfismo}: \omterm{def:b3:topology:continuity}{continuidade} componente a componente nos dois sentidos), de modo que $C
\times C \cong C$.
\end{solution}

\begin{solution}{exo:b3:topology:11}
Para $y \in \R^n$, ponha
\[
\tau(y) = \Bigl(\frac{2y}{\norm y^2 + 1},\ \frac{\norm y^2 -
1}{\norm y^2 + 1}\Bigr) \in \R^n\times\R = \R^{n+1}.
\]
Verifica-se que $\norm{\tau(y)} = 1$, $\tau(y) \neq N$, e
$\sigma(\tau(y)) = y$, $\tau(\sigma(x)) = x$: $\sigma$ e
$\tau$ são mutuamente inversas, ambas \omterm{def:b3:topology:continuity}{contínuas} (fórmulas racionais
com denominadores que não se anulam): $S^n\setminus\{N\} \cong
\R^n$. Estenda a $\hat\sigma \colon S^n \to \widehat{\R^n}$ por
$N \mapsto \infty$: uma bijeção, \omterm{def:b3:topology:continuity}{contínua} em todo ponto de
$S^n\setminus\{N\}$, e em $N$: uma \omterm{def:b3:topology:topology}{vizinhança} básica de
$\infty$ é $\widehat{\R^n}\setminus K$, $K$ \omterm{def:b3:topology:compact}{compacto}, $K
\subseteq \bar B(0, R)$; como $\norm{\sigma(x)}^2 = \frac{1 +
x_{n+1}}{1 - x_{n+1}} \to \infty$ quando $x_{n+1} \to 1$, o conjunto
$\{x \in S^n : x_{n+1} > 1 - \eta\}$ é levado para fora de $\bar B(0,R)$
para $\eta$ pequeno: \omterm{def:b3:topology:continuity}{continuidade}. Uma bijeção \omterm{def:b3:topology:continuity}{contínua} do
\omterm{def:b3:topology:compact}{compacto} $S^n$ no \omterm{def:b3:topology:hausdorff}{espaço de Hausdorff} $\widehat{\R^n}$ é um
\omterm{def:b3:topology:continuity}{homeomorfismo}. Retirando outro ponto $p$: uma rotação de $S^n$
leva $p$ a $N$ (rotações são \omterm{def:b3:topology:continuity}{homeomorfismos}), reduzindo ao
caso calculado: $S^n \setminus \{p\} \cong \R^n$.
\end{solution}

\begin{solution}{exo:b3:topology:12}
(a) $T$ é limitado e fechado: um limite de pontos de $T$
com abscissas $\to x_0 > 0$ permanece no gráfico (localmente fechado)
pela \omterm{def:b3:topology:continuity}{continuidade} de $\sin\frac1x$ em $\intoc01$; um limite
com abscissas $\to 0$ tem ordenada em $\intcc{-1}1$, logo
está no segmento. \omterm{def:b3:topology:compact}{Compacto} por Heine--Borel
(\cref{cor:b3:topology:heineborel}). \omterm{def:b3:topology:interior}{Fecho} do gráfico
$\Gamma$: todo ponto $(0, y)$, $y \in \intcc{-1}1$, é
limite de pontos do gráfico --- resolva $\sin\frac1x = y$ perto de $0$
(a função varre $\intcc{-1}1$ em cada intervalo
$[\frac1{(k+1)\pi}, \frac1{k\pi}]$): $\bar\Gamma = T$.

(b) $\Gamma$ é a imagem \omterm{def:b3:topology:continuity}{contínua} do \omterm{def:b3:topology:connected}{conexo}
$\intoc01$ por $x \mapsto (x, \sin\frac1x)$: \omterm{def:b3:topology:connected}{conexo};
e o \omterm{def:b3:topology:interior}{fecho} de um conjunto \omterm{def:b3:topology:connected}{conexo} é \omterm{def:b3:topology:connected}{conexo}
(\cref{thm:b3:topology:connectedbasics}): $T = \bar\Gamma$
é \omterm{def:b3:topology:connected}{conexo}.

(c) Suponha $\gamma\colon\intcc01\to T$ \omterm{def:b3:topology:continuity}{contínua} com
$\gamma(0) = (0,0)$, $\gamma(1) = (1, \sin1)$, e seja $t^*
= \sup\{t : \gamma_1(t) = 0\}$: por \omterm{def:b3:topology:continuity}{continuidade}, $\gamma_1(t^*)
= 0$ e $\gamma_1 > 0$ em $\intoc{t^*}1$. Para todo
$\delta > 0$, o intervalo $\intoc{t^*}{t^*+\delta}$ tem
$\gamma_1$ assumindo todos os valores de algum $\intoc0\eta$
(teorema do valor intermediário, $\gamma_1(t^* + \delta) > 0$);
em particular, ele contém abscissas da forma
$\frac1{\pi/2 + 2k\pi}$ e $\frac1{-\pi/2 + 2k\pi}$ para
$k$ arbitrariamente grandes, nas quais $\gamma_2 = \pm1$. Logo, em
toda \omterm{def:b3:topology:topology}{vizinhança} à direita de $t^*$, $\gamma_2$ assume ambos os
valores $1$ e $-1$: $\gamma_2$ não tem limite em $t^{*+}$,
contradizendo a \omterm{def:b3:topology:continuity}{continuidade}. Não existe \omterm{def:b3:topology:pathconnected}{caminho}.

(d) $T$ é \omterm{def:b3:topology:connected}{conexo}, mas não é \omterm{def:b3:topology:pathconnected}{conexo por caminhos}: as duas noções
diferem. Para $U \subseteq \R^n$ aberto \omterm{def:b3:topology:connected}{conexo} e $a \in U$,
seja $V$ o conjunto dos pontos de $U$ ligáveis a $a$ por um
\omterm{def:b3:topology:pathconnected}{caminho} em $U$. $V$ é aberto: em torno de $x \in V$, uma bola $B(x, r)
\subseteq U$ é estrelada, e concatenar o \omterm{def:b3:topology:pathconnected}{caminho} até $x$
com um segmento alcança todo ponto da bola. $V$ é
fechado em $U$: se $x \in U \setminus V$, o mesmo argumento da bola
mostra que $B(x, r) \cap V = \varnothing$ (um ponto de
$V$ na bola ligaria $a$ a $x$). Não vazio ($a \in
V$), aberto e fechado no \omterm{def:b3:topology:connected}{conexo} $U$: $V = U$. A curva
do seno escapa disso por ser fechada de \omterm{def:b3:topology:interior}{interior} vazio: seu
ponto ``ruim'' $(0,0)$ não tem bola alguma dentro de $T$ para
transpor as oscilações.
\end{solution}

\begin{solution}{pb:b3:topology:1}
\textbf{1.} Suponha que as expansões de $x, y \in C$ difiram pela primeira vez
na ordem $k \leq m$, digamos $b_k(x) = 1$, $b_k(y) = 0$. Então
\[
x - y \;\geq\; \frac{2}{3^k} - \sum_{j > k}\frac{2}{3^j}
= \frac{2}{3^k} - \frac{1}{3^k} = \frac1{3^k} \geq \frac1{3^m}:
\]
contrapositiva: $\abs{x - y} < 3^{-m}$ força a coincidência até a
ordem $m$.

\textbf{2.} Pela questão 1, $b_n$ é constante em $C \cap (x -
3^{-n}, x + 3^{-n})$: localmente constante, logo \omterm{def:b3:topology:continuity}{contínua}. A
aplicação $x \mapsto (b_n(x))_n \in \{0,1\}^\N$ é \omterm{def:b3:topology:continuity}{contínua}
(componente a componente, \cref{prop:b3:topology:universal}(a)) e
bijetora (unicidade das expansões em dígitos $\{0,2\}$); do \omterm{def:b3:topology:compact}{compacto}
$C$ em um \omterm{def:b3:topology:hausdorff}{espaço de Hausdorff}, é um \omterm{def:b3:topology:continuity}{homeomorfismo}
(\cref{cor:b3:topology:compacthomeo}).

\textbf{3.} \omterm{def:b3:topology:continuity}{Continuidade}: se $\abs{x - y} < 3^{-m}$, os
$m$ primeiros dígitos coincidem, de modo que $\abs{g(x) - g(y)} \leq \sum_{n > m}2^{-n}
= 2^{-m}$. Sobrejetividade: todo $t \in [0,1]$ tem uma expansão
binária $t = \sum b_n2^{-n}$, e $t = g\bigl(\sum
2b_n3^{-n}\bigr)$. Não injetiva: $g$ identifica os dois pontos de Cantor
de dígitos $(0,1,1,1,\dots)$ e $(1,0,0,\dots)$ ---
ambos vão para $\frac12$ (a ambiguidade diádica $0.0111\ldots_2 =
0.1000\ldots_2$).

\textbf{4.} Cada componente de $\Phi$ é \omterm{def:b3:topology:continuity}{contínua} pela mesma
estimativa (seus dígitos são uma subsequência dos $b_n$).
Sobrejetividade: dado $(s, t) \in [0,1]^2$, escolha dígitos binários
$(c_n)$ de $s$ e $(d_n)$ de $t$, e intercale: o ponto $x
\in C$ com $b_{2n-1}(x) = c_n$, $b_{2n}(x) = d_n$ satisfaz
$\Phi(x) = (s,t)$.

\textbf{5.} As lacunas são duas a duas disjuntas, com extremos em
$C$, de modo que a fórmula define $f$ sem ambiguidade em $[0,1]$; nos
extremos de uma lacuna, a fórmula afim devolve $\Phi(u)$,
$\Phi(v)$: coerente com $f = \Phi$ em $C$. Sobrejetividade:
já $f(C) = \Phi(C) = [0,1]^2$.

\textbf{6.} Em $t_0 \notin C$: $t_0$ está numa lacuna aberta em
que $f$ é afim: \omterm{def:b3:topology:continuity}{contínua}. Em $x \in C$: registremos primeiro
duas estimativas.

\emph{(i) Em $C$}: se $\abs{y - x} \leq 3^{-M}$, $y \in C$,
então os dígitos coincidem até $M$, de modo que cada componente de $\Phi(y)
- \Phi(x)$ é no máximo $\sum_{n > \lfloor M/2\rfloor} 2^{-n} =
2^{-\lfloor M/2\rfloor}$.

\emph{(ii) Através de uma lacuna}: uma lacuna $(u,v)$ retirada na etapa $k$
tem comprimento $3^{-k}$, e seus extremos têm dígitos que coincidem até
a ordem $k - 1$ (eles diferem a partir da ordem $k$), de modo que
$\norm{\Phi(v) - \Phi(u)}_\infty \leq
2^{-\lfloor (k-1)/2\rfloor}$.

Agora sejam $\abs{t - x} < 3^{-M}$ com $t$ numa lacuna $(u, v)$ da
etapa $k$, e digamos que $u$ seja o extremo do lado de $x$, de modo que
$\abs{u - x} < 3^{-M}$ e $0 < t - u < 3^{-M}$. Então
\[
\norm{f(t) - \Phi(x)}_\infty
\leq \norm{\Phi(u) - \Phi(x)}_\infty
+ \frac{t - u}{v - u}\,\norm{\Phi(v) - \Phi(u)}_\infty
\leq 2^{-\lfloor M/2\rfloor} + \min\bigl(1,\ 3^{k - M}\bigr)\,
2^{-\lfloor(k-1)/2\rfloor} .
\]
Para $k \geq M$, o último termo é $\leq 2^{-\lfloor(M-1)/2
\rfloor}$; para $k < M$, como $2^{-\lfloor(k-1)/2\rfloor} \leq
\sqrt2\,\cdot 2^{-k/2}$,
\[
3^{k-M}\,2^{-\lfloor(k-1)/2\rfloor}
\leq \sqrt2\,\Bigl(\frac{3}{\sqrt2}\Bigr)^{k} 3^{-M}
\leq \sqrt2\,\Bigl(\frac{3}{\sqrt2}\Bigr)^{M} 3^{-M}
= \sqrt2\; 2^{-M/2}
\]
(a quantidade do meio cresce com $k$). Em todos os casos
$\norm{f(t) - \Phi(x)}_\infty \leq C\,2^{-M/2}$
com uma constante absoluta $C$: fazendo $M \to \infty$, obtém-se a
\omterm{def:b3:topology:continuity}{continuidade} em $x$ (os pontos $t \in C$ estão cobertos por (i)). Logo
$f$ é uma sobrejeção \omterm{def:b3:topology:continuity}{contínua} $[0,1] \to [0,1]^2$ --- de fato,
as estimativas mostram que ela é hölderiana de expoente do tipo $\log 2/(2\log
3)$, mas \omterm{def:b3:topology:continuity}{continuidade} é tudo o que afirmamos.

\textbf{7.} Para $[0,1]^n$: reparta os dígitos de $x \in C$ em
$n$ subsequências intercaladas e repita as questões 4--6
literalmente. Para $\R \to \R^2$: seja $h \colon [0,1] \to [-1,1]^2$
uma sobrejeção \omterm{def:b3:topology:continuity}{contínua} (reescale $f$). Defina $F$ em $[m, m
+ \frac12]$ ($m \in \Z$) como uma cópia de $h$ reescalada para preencher
$[-(m+1), m+1]^2$ para $m \geq 0$ (e simetricamente para $m <
0$), e em $[m + \frac12, m + 1]$ como o segmento afim
que liga os valores nos extremos: $F$ é \omterm{def:b3:topology:continuity}{contínua} (colagem em
peças fechadas, \cref{prop:b3:topology:gluing}(b)) e sua imagem
contém $\bigcup_m [-(m+1), m+1]^2 = \R^2$.

\textbf{8.} $[0,1]$ é \omterm{def:b3:topology:compact}{compacto} e $[0,1]^2$ é \omterm{def:b3:topology:hausdorff}{de Hausdorff}: uma
injeção \omterm{def:b3:topology:continuity}{contínua} é um \omterm{def:b3:topology:continuity}{homeomorfismo} sobre sua imagem
(\cref{cor:b3:topology:compacthomeo} aplicado à correstrição).

\textbf{9.} Retire $\frac12$: $[0,1]\setminus\{\frac12\}$ é
desconexo. Se $h \colon [0,1] \to [0,1]^2$ fosse um
\omterm{def:b3:topology:continuity}{homeomorfismo}, $[0,1]^2\setminus\{h(\frac12)\}$ também seria
desconexo (imagens \omterm{def:b3:topology:continuity}{homeomorfas} de espaços desconexos
são desconexas); mas o quadrado menos um ponto é
\omterm{def:b3:topology:pathconnected}{conexo por caminhos}: ligue dois pontos por um \omterm{def:b3:topology:pathconnected}{caminho} de dois segmentos que evite
o furo. Contradição: $[0,1] \not\cong [0,1]^2$.

\textbf{10.} Pela questão 8, uma sobrejeção \omterm{def:b3:topology:continuity}{contínua} injetora
$[0,1] \to [0,1]^2$ seria um \omterm{def:b3:topology:continuity}{homeomorfismo}, contradizendo a
questão 9. A \omterm{def:b3:topology:compact}{compacidade} entrou exatamente no~\cref{cor:b3:topology:compacthomeo}: é ela que torna \omterm{def:b3:topology:continuity}{contínua} a inversa da
bijeção \omterm{def:b3:topology:continuity}{contínua} (imagens de fechados são compactas, logo
fechadas). Sem \omterm{def:b3:topology:compact}{compacidade}, a conclusão realmente falha: $[0, 2\pi) \to S^1$ é uma bijeção \omterm{def:b3:topology:continuity}{contínua}
que não é um \omterm{def:b3:topology:continuity}{homeomorfismo}.

\textbf{11.} Demonstrado aqui: \emph{existe uma sobrejeção
\omterm{def:b3:topology:continuity}{contínua}, mas nenhuma bijeção \omterm{def:b3:topology:continuity}{contínua}, de $[0,1]$ sobre
$[0,1]^2$; em particular $[0,1] \not\cong [0,1]^2$.} Assim,
a ``dimensão'' não é preservada por sobrejeções \omterm{def:b3:topology:continuity}{contínuas} ---
a cardinalidade e mesmo a \omterm{def:b3:topology:continuity}{continuidade} não a enxergam ---, mas ela
\emph{é} um invariante \omterm{def:b3:topology:topology}{topológico} no nível dos
\omterm{def:b3:topology:continuity}{homeomorfismos}, ao menos entre as dimensões $1$ e $2$, em que a
\omterm{def:b3:topology:connected}{conexidade} após a remoção de um ponto basta. O enunciado geral
--- a \emph{invariância do domínio}: $\R^m \cong \R^n$ implica $m =
n$, e uma injeção \omterm{def:b3:topology:continuity}{contínua} $\R^n \to \R^n$ é aberta ---
é verdadeiro, mas exige \omterm{def:b3:topology:topology}{topologia} \omterm{def:b3:galois:algebraic}{algébrica} (homologia), fora do
alcance deste curso.

\textbf{12.} Em \omterm{def:b3:topology:basis}{base} três, $\frac C2 = \{\sum_na_n3^{-n} : a_n
\in \{0, 1\}\}$ (dividir por dois os dígitos $\{0,2\}$ dá os dígitos
$\{0,1\}$). Para $t \in [0,1]$, tome uma expansão ternária qualquer $t =
\sum_nd_n3^{-n}$, $d_n \in \{0,1,2\}$, e reparta cada dígito
como $d_n = a_n + b_n$ com $a_n, b_n \in \{0,1\}$ ($0 = 0{+}0$,
$1 = 1{+}0$, $2 = 1{+}1$): então $t \in \frac C2 + \frac C2$.
O ponto crucial é que cada dígito se reparte \emph{dentro de} $\{0,1\}$, de modo que
nenhum transporte se propaga e os dígitos podem ser tratados
independentemente. A inclusão recíproca é clara (as somas de dígitos permanecem
$\leq 2$). Reescalando por $2$: $C + C = [0,2]$.

\textbf{13.} A aplicação $(x, y) \mapsto x + y$ leva $C \times C$
sobre $C + C = [0,2]$: a poeira, que não contém nenhum quadrado
(\omterm{def:b3:topology:interior}{interior} vazio: \cref{exo:b3:topology:10}), se projeta ao longo
da antidiagonal sobre um segmento inteiro. Autossemelhança: o
primeiro dígito ternário de um ponto de Cantor é $0$ ou $2$, e
retirá-lo dá $C = \frac C3 \cup \bigl(\frac C3 +
\frac23\bigr)$ --- a equação de ponto fixo que gera
toda figura de $C$.

\textbf{14.} Através de $C \cong \{0,1\}^\N$ (questão 2), a
intercalação $(\varepsilon, \varepsilon') \mapsto
(\varepsilon_1, \varepsilon_1', \varepsilon_2,
\varepsilon_2', \dots)$ é uma bijeção
$\{0,1\}^\N\times\{0,1\}^\N \to \{0,1\}^\N$, \omterm{def:b3:topology:continuity}{contínua} nos dois
sentidos (cada coordenada de saída depende de uma única coordenada
de entrada; \omterm{def:b3:topology:constructions}{topologia produto}). Logo $C\times C \cong C$ e,
indutivamente, $C^n \cong C$. Espaços familiares não compartilham isso:
$\R^2 \not\cong \R$ e $[0,1]^2 \not\cong [0,1]$
(questão 11); produtos infinitos, sim: $([0,1]^\N)^2 \cong
[0,1]^\N$ pela mesma intercalação.

\textbf{15.} $\sum_{k\geq1}2\cdot3^{-2k} = 2\cdot\frac{1/9}{1
- 1/9} = \frac14$: a expansão ternária de $\frac14$ é
$0.\overline{02}$, com dígitos em $\{0,2\}$, de modo que $\frac14 \in
C$; não sendo nem uma expansão finita nem uma de cauda $2$
a partir de certa ordem, ele não é extremo de nenhum terço médio retirado. Os extremos
formam um conjunto enumerável (dois por intervalo retirado, uma infinidade enumerável de
intervalos), ao passo que $C \cong \{0,1\}^\N$ não é enumerável
(argumento diagonal, ou questão 19): quase todo ponto de Cantor
é, como $\frac14$, invisível na figura habitual dos extremos.

\textbf{16.} $g(x) = g(y)$ significa que as sequências binárias
$(b_n(x))$, $(b_n(y))$ representam o mesmo real. Sequências
binárias distintas que representam o mesmo número ocorrem exatamente na
ambiguidade diádica $0.w01^\infty = 0.w10^\infty$: no máximo
duas pré-imagens, e exatamente duas precisamente nos racionais
diádicos de $(0,1)$. Assim $g$ colapsa $C$ sobre $[0,1]$
colando uma infinidade enumerável de pares --- o esqueleto combinatório da
escada de Cantor do~\cref{ch:b3:measure}.

\textbf{17.} $x$ não é isolado em $P$, de modo que $B(x, r/2)$
contém algum $y_0 \in P \setminus \{x\}$; $y_0$ tampouco é
isolado, de modo que $B(x, r/2)$ contém um segundo ponto $y_1
\in P$. Qualquer $\rho \leq \min\bigl(\frac{d(y_0,y_1)}3,
\frac r4\bigr)$ torna $\bar B(y_0, \rho)$ e $\bar B(y_1,
\rho)$ disjuntas e contidas em $B(x, r)$. Ambas estão centradas
em pontos de $P$, onde a mesma extração de dois pontos pode ser
repetida: a perfeição é o suprimento inesgotável de pontos próximos
que mantém a recursão viva para sempre.

\textbf{18.} Construa $F_w$ por indução sobre o comprimento da palavra:
$F_\varnothing = P$, e, dentro de cada bola $B(y_w, \rho_w)$,
a questão 17 fornece duas bolas fechadas disjuntas de raio $\leq
\min(\rho_w/2,\ 2^{-\abs w-1}\operatorname{diam}P)$ centradas
em pontos de $P$; ponha $F_{wi} = \bar B(y_{wi}, \rho_{wi})
\cap P$: não vazia (seu \omterm{ex:b3:groups:actions}{centro}) e compacta. Para $\varepsilon$ infinita, os \omterm{def:b3:topology:compact}{compactos} não vazios encaixados
$F_{\varepsilon_1\cdots\varepsilon_m}$ têm interseção não vazia
(propriedade da interseção finita no \omterm{def:b3:topology:compact}{compacto}
$P$), de diâmetro $0$: um único ponto $h(\varepsilon)$.

\textbf{19.} Palavras que diferem pela primeira vez na ordem $m$ mandam suas
imagens para os dois conjuntos \emph{disjuntos} $F_{w0}, F_{w1}$
(prefixo comum $w$): $h$ é injetiva. Se duas palavras coincidem até a
ordem $m$, ambas as imagens estão num conjunto de diâmetro $\leq
2^{-m}\operatorname{diam}P$: $h$ é \omterm{def:b3:topology:continuity}{contínua}. Logo o não enumerável
$\{0,1\}^\N$ se injeta em $P$: todo espaço métrico
\omterm{def:b3:topology:compact}{compacto} \omterm{prop:b3:galois:perfect}{perfeito} é não enumerável --- $[0,1]$ e $C$
entre eles.

\textbf{20.} $h$ é uma injeção \omterm{def:b3:topology:continuity}{contínua} do \omterm{def:b3:topology:compact}{compacto}
$\{0,1\}^\N$ no \omterm{def:b3:topology:hausdorff}{espaço de Hausdorff} (métrico) $P$:
o~\cref{cor:b3:topology:compacthomeo} a promove a um
\omterm{def:b3:topology:continuity}{homeomorfismo} sobre sua imagem; e $\{0,1\}^\N \cong C$
(questão 2). Todo espaço métrico \omterm{def:b3:topology:compact}{compacto} \omterm{prop:b3:galois:perfect}{perfeito} contém uma
cópia do conjunto de Cantor: a perfeição tem um germe universal, e
ele é o de Cantor.

\textbf{21.} Um espaço métrico \omterm{def:b3:topology:compact}{compacto} enumerável não pode ser
\omterm{prop:b3:galois:perfect}{perfeito} (questão 19), de modo que ele tem um ponto isolado. Em $K =
\{0\}\cup\{\frac1n : n \geq 1\}$, todo ponto $\frac1n$ é
isolado e $0$ não é: pontos isolados podem até ser \omterm{def:b3:topology:interior}{densos}
sem que o espaço seja discreto --- a \omterm{def:b3:topology:compact}{compacidade} mantém seu
limite dentro.

\textbf{22.} Seja $P$ o conjunto dos pontos de condensação de $F$
e $\mathcal B$ a família enumerável dos intervalos abertos de
extremos racionais. Todo ponto de $F$ que não é de condensação está
em algum $I \in \mathcal B$ com $I \cap F$ enumerável, de modo que $F
\setminus P$ está contido na união desses traços enumeráveis, em
quantidade enumerável: enumerável. Como $F$ é não enumerável,
$P \neq \varnothing$; $P$ é fechado (se $x \notin P$, o
intervalo testemunha $I$ não contém ponto de condensação
algum: $I \cap F$ enumerável) e $P \subseteq F$ ($F$ fechado:
um ponto de condensação é, em particular, aderente). $P$ é
\omterm{prop:b3:galois:perfect}{perfeito}: se $I \cap P = \{x\}$ para algum intervalo $I$, então
$I \cap F \subseteq \{x\} \cup (F \setminus P)$ seria
enumerável, contradizendo $x \in P$. Agora a construção em árvore
das questões 17--19 vale literalmente dentro de $P$: as peças
$\bar B(y, \rho) \cap P$ são compactas (subconjuntos fechados e limitados
de $\R$), todo \omterm{ex:b3:groups:actions}{centro} é não isolado em $P$, e o
argumento nunca usou mais do que isso. Logo $\abs P \geq
\abs{\{0,1\}^\N} = \mathfrak c$, e $\abs F \leq \mathfrak
c$ trivialmente: um fechado não enumerável $F \subseteq \R$ tem
cardinalidade exatamente $\mathfrak c$. Conjuntos fechados não podem
testemunhar uma falha da hipótese do \omterm{def:b3:topology:continuity}{contínuo}.

\textbf{23.} Sejam $x \neq y$ em $C$ e escolha $m \geq 0$ com
$3^{-(m+1)} \leq \abs{x - y} < 3^{-m}$. Pela questão 1, os
dígitos $b_n(x) = b_n(y)$ coincidem para $n \leq m$; a primeira
coordenada de $\Phi$ usa $b_1, b_3, \dots$, a segunda $b_2,
b_4, \dots$, de modo que, em cada coordenada, os dois pontos partilham ao
menos $\lfloor m/2\rfloor$ dígitos binários iniciais:
\[
\norm{\Phi(x) - \Phi(y)}_\infty
\leq \sum_{k > \lfloor m/2\rfloor}2^{-k}
= 2^{-\lfloor m/2\rfloor} \leq \sqrt2\cdot2^{-m/2}
= \sqrt2\,\bigl(3^{-m}\bigr)^\alpha
\leq \sqrt2\,\bigl(3\abs{x-y}\bigr)^\alpha,
\]
usando $2^{-m/2} = 3^{-m\alpha}$ (tome logaritmos); e
$\sqrt2\cdot3^\alpha = \sqrt2\cdot\sqrt2 = 2$: a cota
$2\abs{x-y}^\alpha$ em $C$. Mesma lacuna $(u, v)$: $f$ é afim
ali, de modo que, para $t, s \in [u, v]$,
\[
\norm{f(t) - f(s)} = \frac{\abs{t-s}}{v-u}\,
\norm{\Phi(v) - \Phi(u)}
\leq \frac{\abs{t-s}}{v-u}\,2(v-u)^\alpha
\leq 2\abs{t-s}^\alpha,
\]
pois $\frac{\abs{t-s}}{v-u} \leq \bigl(\frac{\abs{t-s}}
{v-u}\bigr)^\alpha$ quando $\abs{t-s} \leq v - u$. Caso geral $t <
s$: se $[t, s]$ encontra $C$, sejam $u'$ e $v'$ o menor
e o maior ponto do \omterm{def:b3:topology:compact}{compacto} $C \cap [t, s]$; então $t$
está numa lacuna (ou num ponto de $C$) cujo extremo direito é
$u'$, $s$ numa cujo extremo esquerdo é $v'$, e
\[
\norm{f(t) - f(s)} \leq 2\abs{t - u'}^\alpha
+ 2\abs{v' - u'}^\alpha + 2\abs{s - v'}^\alpha
\leq 6\abs{t - s}^\alpha ;
\]
se $[t, s]$ não encontra $C$, aplica-se o caso da mesma lacuna. Logo $f$ é
$\alpha$-hölderiana com $\alpha = \frac{\ln2}{2\ln3}$.

\textbf{24.} Suponha $\norm{F(t) - F(s)} \leq
C\abs{t-s}^\alpha$ com $F$ sobrejetora e $\alpha >
\frac12$. Corte $[0,1]$ em $N$ intervalos $I_1, \dots, I_N$ de
comprimento $\frac1N$: cada imagem $F(I_j)$ tem diâmetro $\leq
CN^{-\alpha}$. Dois pontos distintos da grade $G_M =
\bigl\{\bigl(\frac iM, \frac jM\bigr)\bigr\}_{0 \leq i, j
\leq M}$ estão à distância $\geq \frac1M$, de modo que um conjunto de diâmetro
$< \frac1M$ contém no máximo um deles. Escolha $M =
\lfloor N^\alpha/(2C)\rfloor$ (com $N$ grande): então
$CN^{-\alpha} \leq \frac1{2M} < \frac1M$, e a sobrejetividade
coloca cada um dos $(M+1)^2$ pontos da grade em algum
$F(I_j)$, donde
\[
\frac{N^{2\alpha}}{4C^2} \leq (M + 1)^2 \leq N,
\]
e $N^{2\alpha - 1} \leq 4C^2$ falha para $N$ grande quando
$2\alpha > 1$: contradição. Nossa curva, com $\alpha \approx
0.3155$, fica abaixo do muro, como deve ser; a curva de Hilbert
(admitida) é $\frac12$-hölderiana, de modo que o expoente crítico
$\frac12$ é atingido --- a regularidade hölderiana, ao contrário da
injetividade, é uma questão de grau, e $\frac12$ é exatamente
a fronteira que a dimensão $2$ impõe a uma aplicação vinda da
dimensão $1$.

\textbf{25.} O $m$-ésimo aglomerado $\{\frac1m + \frac1n : n >
m^2\}$ está em $\bigl(\frac1m, \frac1m + \frac1{m^2}\bigr]
\subseteq \bigl(\frac1m, \frac1{m-1}\bigr)$ para $m \geq 2$,
pois $\frac1{m-1} - \frac1m = \frac1{m(m-1)} >
\frac1{m^2}$ (e o primeiro aglomerado está em $(1, \frac32]$):
os aglomerados ocupam intervalos disjuntos. Pontos limites de $K$:
dentro do $m$-ésimo intervalo, apenas $\frac1m$ (o aglomerado
converge para ele e é discreto em si mesmo); globalmente, $0$
(toda \omterm{def:b3:topology:topology}{vizinhança} de $0$ contém aglomerados inteiros). Assim $K'
= \{0\} \cup \{\frac1m\}$, depois $K'' = \{0\}$ (cada $\frac1m$
é isolado em $K'$), $K''' = \varnothing$; $K$ contém seus
pontos limites, logo é fechado, limitado, enumerável, \omterm{def:b3:topology:compact}{compacto},
e seus pontos isolados --- os pontos dos aglomerados $\frac1m +
\frac1n$ --- são \omterm{def:b3:topology:interior}{densos} em $K$: todo elemento de $K'$ é
limite deles, como o mecanismo da questão 21 prevê. Indução:
$K_1 = \{0\}\cup\{\frac1n\}$ tem $K_1^{(1)} = \{0\}$; dado
um \omterm{def:b3:topology:compact}{compacto} enumerável $K_j \subseteq [0, 1]$ com $K_j^{(j)} =
\{0\}$ e $K_j^{(j+1)} = \varnothing$, ponha
\[
K_{j+1} = \{0\} \cup \bigcup_{m \geq 1}
\Bigl(\frac1m + \varepsilon_m K_j\Bigr),
\]
com $\varepsilon_m$ pequeno o bastante para que a $m$-ésima cópia esteja
em $\bigl(\frac1m, \frac1{m-1}\bigr)$. A derivação age cópia
a cópia (as cópias vivem em intervalos abertos disjuntos), de modo que
$K_{j+1}^{(i)} = \{0\} \cup \bigcup_m\bigl(\frac1m +
\varepsilon_mK_j^{(i)}\bigr)$ para $i \leq j$; em $i = j$ isso
se lê $\{0\} \cup \{\frac1m\}$, e então $K_{j+1}^{(j+1)} =
\{0\}$ e $K_{j+1}^{(j+2)} = \varnothing$. Toda ordem finita
ocorre. Um conjunto \omterm{prop:b3:galois:perfect}{perfeito} é o outro extremo: $P' = P$, a
derivação nunca se move --- e Cantor--Bendixson (questão
22) diz precisamente que todo fechado se reparte em um
núcleo \omterm{prop:b3:galois:perfect}{perfeito}, invisível à derivação, e um resto enumerável
que a derivação vai comendo.
\end{solution}
